\section{最小范数插值}
\label{sec:rkhs-min-norm-interpolation}

现在回到第~\ref{sec:rkhs-recovery}节的精确插值问题，把候选空间取为RKHS。固定$X$上的RKHS $H$及其再生核$K$，取观测位置$x_1,\ldots,x_n\in X$（暂不要求互异）和读数$\bm y\in\mathbb K^n$，要找$f\in H$使$f(x_i)=y_i$对每个$i$成立。

这个问题要分两步回答。第一步判断数据能否由$H$中的函数实现；若不能，谈论“最好的插值”没有意义。第二步在可行函数中选出范数最小的一个。两步都依赖同一个正交分解：把$H$分成读数能看见的有限维部分与读数完全看不见的部分。

\subsection{观测映射}
\label{sec:rkhs-observation-map}

把$n$个读数合在一起，得到线性映射
\[
S:H\to\mathbb K^n,\qquad Sf=\bigl(f(x_1),\ldots,f(x_n)\bigr)^{\top},
\]
称为\term{观测映射}（observation map）。精确插值就是解方程$Sf=\bm y$。由再生性质，$Sf$的第$i$个分量是$\langle K_{x_i},f\rangle$，所以$S$是有界线性映射，并且$\lVert Sf\rVert_2^2\leq\bigl(\sum_iK(x_i,x_i)\bigr)\lVert f\rVert_H^2$。

\paragraph{伴随与Gram矩阵}

$\mathbb K^n$取标准内积$\langle\bm c,\bm d\rangle=\bm c^*\bm d$。对任意$\bm c\in\mathbb K^n$，
\[
\langle\bm c,Sf\rangle=\sum_i\overline{c_i}\,f(x_i)=\sum_i\overline{c_i}\,\langle K_{x_i},f\rangle_H=\Bigl\langle\sum_ic_iK_{x_i},\ f\Bigr\rangle_H ,
\]
所以$S$的伴随把系数映成核截面的组合：$S^*\bm c=\sum_ic_iK_{x_i}$。这类函数的读数和内积都可以用Gram矩阵$\bm G_{i,j}=K(x_i,x_j)$计算：
\begin{equation}
S(S^*\bm c)=\bm G\bm c,\qquad
\langle S^*\bm c,S^*\bm d\rangle_H=\bm c^*\bm G\bm d,\qquad
SS^*=\bm G .
\label{eq:rkhs-gram-calculus}
\end{equation}
第一式的第$i$个分量是$\sum_jc_jK(x_i,x_j)$，第二式由$\langle K_{x_i},K_{x_j}\rangle=K(x_i,x_j)$展开得到。于是在核截面的张成空间内，函数的全部计算都化为$n$维向量与$n\times n$矩阵的计算。注意$\bm c^*\bm G\bm c$才是函数的范数平方，系数的欧氏长度$\lVert\bm c\rVert_2$一般与函数范数无关。

\paragraph{秩与线性相关}

$\bm G$的秩恰好是核截面$K_{x_1},\ldots,K_{x_n}$张成空间的维数。因为$\bm G$是Hermite半正定矩阵，$\bm c^*\bm G\bm c=0$当且仅当$\bm G\bm c=\bm 0$；而由式~\eqref{eq:rkhs-gram-calculus}，$\bm c^*\bm G\bm c=\lVert S^*\bm c\rVert_H^2$。所以$\ker\bm G=\ker S^*$，即$\bm G\bm c=\bm 0$恰好描述核截面之间的线性关系，由秩—零度公式即得上述结论。

$\bm G$奇异有三种典型来源。一是重复观测：若$x_1=x_2$，则$K_{x_1}=K_{x_2}$，$\bm G$的前两行相同。二是核截面所在的空间维数不够：仿射核在$-1,0,1$上的Gram矩阵以$(1,-2,1)^{\top}$为零向量，因为二维空间中三个截面必然线性相关。三是退化的观测点：最小值核满足$K_0=\min(\cdot,0)=0$，$H_{\min}$中所有函数在$0$处取零，含有$x=0$的观测点组，Gram矩阵有一整行为零。

\begin{example}[最小值核的观测映射]
\label{ex:rkhs-hmin-observation}
在$H_{\min}$中取$x_1,x_2,x_3=\tfrac13,\tfrac23,1$，Gram矩阵为
\[
\bm G=\frac13\begin{bmatrix}1&1&1\\1&2&2\\1&2&3\end{bmatrix},\qquad
\det\bm G=\frac1{27}\neq0 .
\]
三个核截面线性无关，它们的组合$S^*\bm c=\sum_jc_j\min(\cdot,x_j)$恰好是左端为零、只在$\tfrac13,\tfrac23$处转折的连续折线：在三段上的斜率依次为$c_1+c_2+c_3$、$c_2+c_3$、$c_3$。
\end{example}

\subsection{观测零空间}
\label{sec:rkhs-observation-null-space}

读数看不见的函数组成$S$的核
\[
N=\ker S=\{f\in H:f(x_1)=\cdots=f(x_n)=0\},
\]
称为\term{观测零空间}（observation null space）。$S$有界，所以$N$是闭子空间。记$M=\operatorname{span}\{K_{x_1},\ldots,K_{x_n}\}$，它是$S^*$的像，有限维因而闭。下面的定理说明这两个子空间互为正交补。

\begin{theorem}[观测零空间的正交补]
\label{thm:rkhs-observation-null-space}
在上述记号下，$N=M^\perp$，$N^\perp=M$，从而$H=M\oplus N$。记$P$为到$M$的正交投影，则对一切$f\in H$，$S(Pf)=Sf$且$\lVert Pf\rVert_H\leq\lVert f\rVert_H$。两个函数$f,g$在全部观测点上读数相同，当且仅当$f-g\in N$，也当且仅当$Pf=Pg$。
\end{theorem}

\begin{proof}
$f\in N$当且仅当$\langle K_{x_i},f\rangle=0$对每个$i$成立，即$f$与$M$的每个生成元正交，这就是$N=M^\perp$。$M$是闭子空间，所以$N^\perp=M^{\perp\perp}=M$，正交分解$H=M\oplus N$由第~\ref{chap:inner-product-spaces}章的投影定理给出。$f-Pf\in N$说明$S(f-Pf)=0$；勾股定理给出$\lVert f\rVert^2=\lVert Pf\rVert^2+\lVert f-Pf\rVert^2$。最后，$Sf=Sg$等价于$f-g\in N$，而由正交分解，这又等价于$P(f-g)=0$。
\end{proof}

所以有限个点值恰好确定投影$Pf$，余项$f-Pf$完全看不见：它在每个观测点取零，不论多大都不改变任何读数。

\begin{example}[折线与观测零空间的正交性]
\label{ex:rkhs-hmin-null-space}
继续例~\ref{ex:rkhs-hmin-observation}。$N$由$H_{\min}$中在$\tfrac13,\tfrac23,1$处取零的函数组成，例如$\sin(3k\pi x)$（$k=1,2,\ldots$），所以$N$是无限维的。$M$是左端为零、以$\tfrac13,\tfrac23$为节点的折线。

正交性可以直接验证：设$h\in M$在第$k$段上的斜率为$s_k$，$g\in N$，则
\[
\langle h,g\rangle_{H_{\min}}=\int_0^1h'g'=\sum_{k=1}^3s_k\bigl(g(\tfrac k3)-g(\tfrac{k-1}3)\bigr)=0 ,
\]
因为$g$在四个节点$0,\tfrac13,\tfrac23,1$上都取零。在这个例子中，$Pf$就是过原点和三个观测点的折线插值。
\end{example}

定理~\ref{thm:rkhs-observation-null-space}还有一个更一般的推论。若某个准则只通过读数$Sf$依赖于$f$，再加上一项随$\lVert f\rVert_H$增大的惩罚，那么把$f$换成$Pf$不改变前者、不增大后者，最优函数总可以取在有限维空间$M$中，即写成$\sum_jc_jK_{x_j}$的形式。这就是带损失与正则项的函数拟合中表示定理的结构依据，可参考\BookChapterRef{chap:function-fitting-regularization}{卷五第10章“函数拟合与正则化”}。本章只讨论精确插值这一特殊情形。

\subsection{约束相容性}
\label{sec:rkhs-consistency}

给定读数$\bm y$，第一步要判断是否存在$f\in H$满足$Sf=\bm y$。不是任意数据都能被插值：在重复观测点上给出两个不同的读数，任何函数都无法同时满足。下面的定理把可行性写成矩阵条件。

\begin{theorem}[插值约束的相容性]
\label{thm:rkhs-consistency}
存在$f\in H$使$Sf=\bm y$，当且仅当$\bm y$属于$\bm G$的像，也当且仅当$\bm d^*\bm y=0$对一切$\bm d\in\ker\bm G$成立。
\end{theorem}

\begin{proof}
若$Sf=\bm y$，由定理~\ref{thm:rkhs-observation-null-space}，$\bm y=S(Pf)$，而$Pf\in M$可写成$S^*\bm c$，于是$\bm y=SS^*\bm c=\bm G\bm c$。反过来，若$\bm y=\bm G\bm c$，则$f=S^*\bm c$满足$Sf=\bm y$。$\bm G$是Hermite矩阵，其像是$\ker\bm G$的正交补，这给出第二个等价条件。
\end{proof}

当$\bm G$可逆时，像是整个$\mathbb K^n$，任意数据都可以插值；特别地，严格正定核在互异观测点上总是如此。当$\bm G$奇异时，数据必须满足$\ker\bm G$给出的线性关系。仿射核在$-1,0,1$上，$\ker\bm G$由$(1,-2,1)^{\top}$张成，相容条件是$y_1-2y_2+y_3=0$，即三个数据点共线。读数$(0,1,0)$不满足它，所以没有仿射函数通过这三点；这与第~\ref{sec:rkhs-finite-feature-space}节例~\ref{ex:rkhs-polynomial-features}的第1种情形是同一现象，只是那里用设计矩阵的列空间来判断。重复观测$x_1=x_2$时$(1,-1,0,\ldots,0)^{\top}\in\ker\bm G$，相容条件是$y_1=y_2$；$H_{\min}$中在$0$处观测时，相容条件是该读数为零。

还要区分“没有插值函数”与“插值函数不止一个”，二者由不同的对象决定。前者取决于$\bm y$是否落在$\bm G$的像中；后者取决于观测零空间$N$是否非零，与$\bm G$是否可逆无关。仿射核在$-1,0,1$上$\bm G$奇异，但三个截面已经张成整个$H_{\mathrm{aff}}$，$N=\{0\}$，相容数据确定唯一的函数。$H_{\min}$中三个观测点给出可逆的$\bm G$，任何数据都可行，$N$却是无限维的，可行函数有无穷多个。所以不能把$\bm G$奇异理解为插值不唯一，也不能把$\bm G$可逆理解为函数被唯一确定。

\subsection{最小范数解}
\label{sec:rkhs-min-norm-solution}

数据相容时，可行函数组成仿射集$f_0+N$，其中$f_0$是任意一个可行函数。第~\ref{sec:rkhs-finite-feature-space}节在有限特征空间中遇到解不唯一时，卷二的做法是选欧氏长度最小的系数；在RKHS中，自然的做法是选范数最小的函数。

\begin{theorem}[最小范数插值]
\label{thm:rkhs-min-norm-interpolation}
设$\bm y$属于$\bm G$的像，$\bm c$是方程$\bm G\bm c=\bm y$的任意一个解，令
\[
f^+=S^*\bm c=\sum_{j=1}^nc_jK_{x_j},\qquad f^+(x)=\sum_{j=1}^nc_jK(x,x_j).
\]
则$f^+$不依赖于解$\bm c$的选取，它是$M$中唯一满足$Sf=\bm y$的函数；对任意满足$Sg=\bm y$的$g\in H$，
\begin{equation}
\lVert g\rVert_H^2=\lVert f^+\rVert_H^2+\lVert g-f^+\rVert_H^2 ,
\label{eq:rkhs-min-norm-pythagoras}
\end{equation}
所以$f^+$是唯一的最小范数可行函数，并且$\lVert f^+\rVert_H^2=\bm c^*\bm y=\bm y^*\bm G^{\dagger}\bm y$，其中$\bm G^{\dagger}$是Moore--Penrose伪逆。
\end{theorem}

\begin{proof}
由式~\eqref{eq:rkhs-gram-calculus}，$Sf^+=\bm G\bm c=\bm y$，所以$f^+$可行。若$\bm c'$是另一个解，则$\bm G(\bm c-\bm c')=\bm 0$，由上一小节$S^*(\bm c-\bm c')=0$，两组系数给出同一函数。$M$中两个读数相同的函数之差属于$M\cap N=\{0\}$，这给出$M$中的唯一性。

对可行的$g$，$g-f^+\in N$，而$f^+\in M=N^\perp$，勾股定理即为式~\eqref{eq:rkhs-min-norm-pythagoras}，等号仅在$g=f^+$时成立。最后，$\lVert f^+\rVert^2=\bm c^*\bm G\bm c=\bm c^*\bm y$；取$\bm c=\bm G^{\dagger}\bm y$是方程的一个解（$\bm y$在$\bm G$的像中），代入并用$\bm G^{\dagger}$为Hermite矩阵，得$\bm y^*\bm G^{\dagger}\bm y$。
\end{proof}

称$f^+$为数据的\term{最小范数插值}（minimum-norm interpolation）。由定理~\ref{thm:rkhs-observation-null-space}，它也是任意可行函数在$M$上的投影。整个求解只需要核在观测点上的值：解一个$n\times n$的线性方程，再用核截面组合出函数。

\paragraph{系数不唯一与函数唯一}

定理区分了两种唯一性。$\bm G$奇异时，方程$\bm G\bm c=\bm y$的解组成$\bm c_0+\ker\bm G$，系数不唯一；但$\ker\bm G$中的系数给出零函数，所以函数$f^+$唯一。伪逆从全部解中挑出欧氏长度最小的系数，这只是选了一种写法，与函数范数最小是两回事。

仿射核在$-1,0,1$上，取相容数据$\bm y=(0,1,2)^{\top}$。$\bm c=(0,0,1)^{\top}$是一个解，全部解为$(s,-2s,1+s)^{\top}$。无论$s$取何值，$f^+=sK_{-1}-2sK_0+(1+s)K_1=1+t$，范数平方$\bm c^*\bm y=-2s+2(1+s)=2=1^2+1^2$。

\begin{example}[最小值核的折线插值]
\label{ex:rkhs-hmin-interpolation}
在例~\ref{ex:rkhs-hmin-observation}的观测点上，取$q(x)=4x(1-x)$的读数$\bm y=(\tfrac89,\tfrac89,0)^{\top}$。由定理~\ref{thm:rkhs-min-norm-interpolation}与例~\ref{ex:rkhs-hmin-null-space}，$f^+$是过原点和三个数据点的折线，三段斜率依次为$\tfrac83,0,-\tfrac83$。对照例~\ref{ex:rkhs-hmin-observation}中的斜率公式，得$\bm c=(\tfrac83,\tfrac83,-\tfrac83)^{\top}$。

代回检验$\bm G\bm c=\bm y$：第一行$\tfrac13\cdot\tfrac83=\tfrac89$，第二行$\tfrac13(\tfrac83+\tfrac{16}3-\tfrac{16}3)=\tfrac89$，第三行$\tfrac13(\tfrac83+\tfrac{16}3-8)=0$。范数为$\lVert f^+\rVert^2=\bm c^{\top}\bm y=\tfrac{128}{27}$，也等于折线导数的平方积分$2\cdot\tfrac13\cdot\tfrac{64}{9}$。而$\lVert q\rVert_{H_{\min}}^2=\int_0^1(4-8x)^2\,\mathrm dx=\tfrac{16}3=\tfrac{144}{27}$，由式~\eqref{eq:rkhs-min-norm-pythagoras}，$\lVert q-f^+\rVert^2=\tfrac{16}{27}$。
\end{example}

最小范数插值只回答了“在可行函数中选哪一个”。它在测点之间与真实函数相差多少，取决于真实函数在$N$中的分量，这是读数看不见的部分。下一节在已知目标函数的前提下研究这种表示误差，并看它怎样随子空间的扩大而变化。
